% ============================================================
%  C: Como Programar -- 6ª Edição | Deitel & Deitel
%  Capítulo 10 -- Estruturas, Uniões, Manipulações de Bits
%               e Enumerações em C
%  EXERCÍCIOS DE AUTORREVISÃO (10.1 -- 10.4)
%  Abordagem incremental: Caps. 1 a 10
% ============================================================
\documentclass[12pt,a4paper]{article}
\usepackage[utf8]{inputenc}
\usepackage[T1]{fontenc}
\usepackage[brazil]{babel}
\usepackage{geometry}
\geometry{top=2.5cm,bottom=2.5cm,left=3cm,right=3cm}
\usepackage{parskip}\setlength{\parskip}{6pt}\setlength{\parindent}{0pt}
\usepackage{xcolor,tcolorbox,enumitem,listings,fancyhdr,titlesec,hyperref,booktabs,array}
\tcbuselibrary{skins,breakable}

\definecolor{deitelblue}{RGB}{0,70,127}
\definecolor{deitelgray}{RGB}{240,242,245}
\definecolor{deitelgreen}{RGB}{0,110,60}
\definecolor{deitellightblue}{RGB}{220,235,250}
\definecolor{deitelred}{RGB}{160,20,20}
\definecolor{deitellorange}{RGB}{200,90,0}

\newtcolorbox{boaprática}[1][]{colback=deitellightblue,colframe=deitelblue,
  fonttitle=\bfseries\small,title={Boa prática de programação},breakable,#1}
\newtcolorbox{erroco}[1][]{colback=red!5,colframe=deitelred,
  fonttitle=\bfseries\small,title={Erro comum de programação},breakable,#1}
\newtcolorbox{respostabox}[1][]{colback=deitelgray,colframe=deitelblue!60,
  fonttitle=\bfseries\small\color{deitelblue},title={Resposta detalhada},breakable,#1}
\newtcolorbox{enunciadoBox}[1][]{colback=white,colframe=deitelblue,
  fonttitle=\bfseries,title={#1},breakable}
\newtcolorbox{saida}[1][]{colback=black!88,colframe=deitelblue,
  fontupper=\ttfamily\small\color{green!70},title={\small\color{white}Saída do programa},breakable,#1}

\setlist[enumerate,1]{label=\textbf{\alph*)},leftmargin=2em}
\setlist[itemize]{leftmargin=2em}

\lstset{
  basicstyle=\ttfamily\small,backgroundcolor=\color{deitelgray},
  frame=single,framerule=0.4pt,rulecolor=\color{deitelblue!40},
  breaklines=true,showstringspaces=false,
  keywordstyle=\color{deitelblue}\bfseries,
  commentstyle=\color{deitelgreen}\itshape,
  stringstyle=\color{deitelred},
  language=C,numbers=left,numberstyle=\tiny\color{gray},
  xleftmargin=18pt,xrightmargin=5pt,stepnumber=1,
}

\pagestyle{fancy}\fancyhf{}
\fancyhead[L]{\small\color{deitelblue}\textbf{C: Como Programar -- 6ª ed. | Deitel \& Deitel}}
\fancyhead[R]{\small\color{deitelblue}Cap. 10 -- Autorrevisão}
\fancyfoot[C]{\small\thepage}
\renewcommand{\headrulewidth}{0.4pt}

\titleformat{\section}[block]{\large\bfseries\color{deitelblue}}{\thesection.}{0.5em}{}[\titlerule]
\titleformat{\subsection}[block]{\normalsize\bfseries\color{deitelblue!80}}{}{}{}

\hypersetup{colorlinks=true,linkcolor=deitelblue}

\begin{document}

\begin{titlepage}
  \centering\vspace*{2cm}
  {\color{deitelblue}\rule{\linewidth}{2pt}}\\[0.4cm]
  {\huge\bfseries\color{deitelblue}C: Como Programar}\\[0.2cm]
  {\large\color{deitelblue}6ª Edição $\cdot$ Paul Deitel \& Harvey Deitel}\\[0.2cm]
  {\color{deitelblue}\rule{\linewidth}{2pt}}\\[1cm]
  {\LARGE\bfseries Capítulo 10}\\[0.3cm]
  {\Large\color{deitelblue!80}Estruturas, Uniões, Manipulações de Bits e Enumerações em C}\\[0.8cm]
  \begin{tcolorbox}[colback=deitellightblue,colframe=deitelblue,width=0.85\linewidth]
    \centering{\large\bfseries Exercícios de Autorrevisão (10.1--10.4)}\\[0.2cm]
    {\normalsize Resolvidos passo a passo, abordagem incremental\\
    Capítulos 1 a 10}
  \end{tcolorbox}
\end{titlepage}

\section*{Construindo sobre os Capítulos Anteriores}

No Capítulo~10, ampliamos drasticamente nossa capacidade de modelar dados em C.
Construindo sobre os fundamentos dos Caps.~1--9, incorporamos quatro famílias de
recursos poderosos:

\begin{itemize}[noitemsep]
  \item \textbf{Estruturas (\texttt{struct})} -- agrupam variáveis \textit{heterogêneas}
  (tipos diferentes) sob um único nome, criando ``registros''
  \item \textbf{Uniões (\texttt{union})} -- estruturas onde os membros
  \textit{compartilham} o mesmo espaço de memória
  \item \textbf{Operadores sobre bits} -- \texttt{\&}, \texttt{|}, \texttt{\^{}},
  \texttt{\~{}}, \texttt{<<}, \texttt{>>} para manipulação em nível de bits
  \item \textbf{Campos de bit} -- membros de estrutura com largura específica em bits
  \item \textbf{\texttt{typedef}} -- cria sinônimos para tipos existentes
  \item \textbf{Enumerações (\texttt{enum})} -- conjuntos nomeados de constantes inteiras
\end{itemize}

\begin{boaprática}
  As \texttt{struct} são a ponte para conceitos de programação orientada a objetos.
  Antes de \texttt{struct}, manipulávamos apenas tipos primitivos isolados ou arrays
  homogêneos. Agora podemos criar \textit{tipos próprios} -- a base do design moderno
  de software.
\end{boaprática}

\tableofcontents\newpage

% ============================================================
\section{Exercício 10.1 -- Preenchimento de Lacunas (14 itens)}
% ============================================================

\begin{respostabox}
\begin{enumerate}[label=\textbf{\alph*)}]

\item Uma \textbf{estrutura} (\texttt{struct}) é uma coleção de variáveis relacionadas
agrupadas sob um único nome.

\item Uma \textbf{união} (\texttt{union}) é uma coleção de variáveis sob um único nome
que \textit{compartilham o mesmo espaço de armazenamento}.

\item Os bits no resultado de uma expressão com o operador \textbf{AND sobre bits
(\texttt{\&})} são 1 se os bits correspondentes em \textit{ambos} os operandos são 1.

\item Variáveis declaradas em uma definição de estrutura são chamadas de \textbf{membros}.

\item Em expressões com \textbf{OR inclusivo sobre bits (\texttt{|})}, os bits são 1
se \textit{pelo menos um} dos bits correspondentes for 1.

\item A palavra-chave \textbf{\texttt{struct}} introduz uma declaração de estrutura.

\item A palavra-chave \textbf{\texttt{typedef}} cria um sinônimo para um tipo previamente
definido.

\item Em expressões com \textbf{OR exclusivo (XOR) sobre bits (\texttt{\^{}})}, os bits são
1 se \textit{exatamente um} dos bits correspondentes for 1.

\item O AND sobre bits é normalmente usado para \textbf{mascarar} bits (selecionar
certos e zerar outros).

\item A palavra-chave \textbf{\texttt{union}} introduz uma definição de união.

\item O nome da estrutura é chamado de \textbf{nome de tag} (\textit{tag name}).

\item Um membro pode ser acessado com o operador \textbf{ponto (\texttt{.})} ou
\textbf{seta (\texttt{->})}, este último quando temos um ponteiro para a estrutura.

\item \textbf{\texttt{<<}} (deslocamento à esquerda) e \textbf{\texttt{>>}} (deslocamento
à direita) deslocam os bits de um valor.

\item Uma \textbf{enumeração} (\texttt{enum}) é um conjunto de inteiros representados
por identificadores.

\end{enumerate}
\end{respostabox}

\newpage

% ============================================================
\section{Exercício 10.2 -- Verdadeiro ou Falso}
% ============================================================

\subsection*{Item (a) -- \texttt{struct} só com um tipo}
\begin{tcolorbox}[colback=red!5,colframe=deitelred,title={\bfseries FALSO}]
Esta é justamente a grande vantagem das estruturas sobre arrays: \texttt{struct}
pode conter \textbf{muitos tipos diferentes}. Esse é o ponto-chave que diferencia
\texttt{struct} de array (que é homogêneo).

\begin{lstlisting}
struct aluno {
    char nome[ 50 ];   /* char[] */
    int  matricula;    /* int    */
    double media;       /* double */
};
\end{lstlisting}
\end{tcolorbox}

\subsection*{Item (b) -- Uniões podem ser comparadas com \texttt{==}}
\begin{tcolorbox}[colback=red!5,colframe=deitelred,title={\bfseries FALSO}]
Uniões \textbf{não} podem ser comparadas diretamente porque podem ter
\textit{bytes indefinidos} (padding) com valores diferentes mesmo quando o
conteúdo lógico é o mesmo. A comparação byte-a-byte poderia falhar incorretamente.

Para comparar uniões, compare membros individuais de forma explícita.
\end{tcolorbox}

\subsection*{Item (c) -- Nome da tag é opcional}
\begin{tcolorbox}[colback=green!5,colframe=deitelgreen,title={\bfseries VERDADEIRO}]
Uma estrutura pode ser declarada sem nome de tag (estrutura anônima), útil
geralmente combinada com \texttt{typedef}:
\begin{lstlisting}
typedef struct {       /* sem tag */
    int x, y;
} Ponto;
\end{lstlisting}
\end{tcolorbox}

\subsection*{Item (d) -- Membros de \texttt{struct}s diferentes devem ter nomes únicos}
\begin{tcolorbox}[colback=red!5,colframe=deitelred,title={\bfseries FALSO}]
Membros de \textit{estruturas diferentes} podem ter nomes iguais sem conflito,
porque o acesso é qualificado pelo nome da estrutura: \texttt{a.nome} vs
\texttt{b.nome}. Apenas dentro da \textbf{mesma} estrutura os nomes devem ser
únicos.

\begin{lstlisting}
struct ponto2d { double x, y; };
struct ponto3d { double x, y, z; };  /* x e y aqui sao OK */
\end{lstlisting}
\end{tcolorbox}

\subsection*{Item (e) -- \texttt{typedef} define novos tipos}
\begin{tcolorbox}[colback=red!5,colframe=deitelred,title={\bfseries FALSO}]
\texttt{typedef} \textit{não cria} novos tipos -- ele apenas cria \textbf{sinônimos}
(\textit{aliases}) para tipos existentes. \texttt{Idade} e \texttt{int} são
intercambiáveis se \texttt{typedef int Idade;}.

\begin{lstlisting}
typedef int Idade;
Idade a = 30;      /* equivale a int a = 30; */
int   b = a;       /* sem conversao -- mesmo tipo */
\end{lstlisting}
\end{tcolorbox}

\subsection*{Item (f) -- \texttt{struct}s sempre passadas por referência}
\begin{tcolorbox}[colback=red!5,colframe=deitelred,title={\bfseries FALSO}]
Estruturas são passadas \textbf{por valor} por padrão em C -- exatamente como
qualquer outro tipo. Isso pode ser \textit{caro} para estruturas grandes, pois
\textit{toda a estrutura é copiada}.

Para passar por referência, passe um ponteiro explicitamente:
\begin{lstlisting}
void f( struct grande s );          /* por valor: COPIA tudo */
void f( struct grande *s );         /* por referencia */
\end{lstlisting}

\begin{boaprática}
  Para \texttt{struct} grandes, prefira passar por ponteiro. Use \texttt{const}
  quando não pretende modificar a estrutura -- princípio do menor privilégio.
\end{boaprática}
\end{tcolorbox}

\subsection*{Item (g) -- \texttt{struct}s não podem ser comparadas com \texttt{==}}
\begin{tcolorbox}[colback=green!5,colframe=deitelgreen,title={\bfseries VERDADEIRO}]
\texttt{struct}s não podem ser comparadas diretamente com \texttt{==} ou
\texttt{!=}, pelos mesmos motivos das uniões: bytes de \textit{padding} (preenchimento
para alinhamento) podem ter valores indefinidos.

Para comparar duas \texttt{struct}, compare membro por membro:
\begin{lstlisting}
if ( a.x == b.x && a.y == b.y ) { ... }
\end{lstlisting}
\end{tcolorbox}

\newpage

% ============================================================
\section{Exercício 10.3 -- Código com \texttt{struct peça}}
% ============================================================

\begin{respostabox}
\begin{enumerate}[label=\textbf{\alph*)}]

\item \textbf{Definir \texttt{struct peça}:}
\begin{lstlisting}
struct peca {
    int  numPeca;
    char nomePeca[ 26 ];  /* 25 chars + '\0' */
};
\end{lstlisting}

\item \textbf{\texttt{typedef} para criar \texttt{Peça}:}
\begin{lstlisting}
typedef struct peca Peca;
\end{lstlisting}

\item \textbf{Declarações usando \texttt{Peça}:}
\begin{lstlisting}
Peca a;          /* variavel simples */
Peca b[ 10 ];    /* array de 10 Pecas */
Peca *ptr;       /* ponteiro para Peca */
\end{lstlisting}
\textit{Nota:} com \texttt{typedef} podemos colocar tudo em uma linha:
\texttt{Peca a, b[10], *ptr;}.

\item \textbf{Ler número e nome para os membros de \texttt{a}:}
\begin{lstlisting}
scanf( "%d%25s", &a.numPeca, a.nomePeca );
\end{lstlisting}
\textit{Cuidado:} \texttt{a.nomePeca} (sem \texttt{\&}) porque é um array;
o nome já é o endereço.

\item \textbf{Atribuir \texttt{a} ao elemento 3 de \texttt{b}:}
\begin{lstlisting}
b[ 3 ] = a;   /* C suporta atribuicao entre structs */
\end{lstlisting}

\item \textbf{Atribuir endereço de \texttt{b} a \texttt{ptr}:}
\begin{lstlisting}
ptr = b;   /* nome do array = endereco do 1o elemento */
\end{lstlisting}

\item \textbf{Imprimir elemento 3 via \texttt{ptr}:}
\begin{lstlisting}
printf( "%d %s\n", ( ptr + 3 )->numPeca,
                   ( ptr + 3 )->nomePeca );
\end{lstlisting}

\textit{Explicação:} \texttt{ptr + 3} é aritmética de ponteiro (Cap.~7) -- aponta
para \texttt{b[3]}. O operador \texttt{->} é equivalente a \texttt{(*p).membro},
mas mais legível.
\end{enumerate}

\begin{boaprática}
  A combinação \texttt{struct} + \texttt{typedef} + ponteiro é o tripé que
  sustenta todo o desenvolvimento orientado a tipos em C. Listas encadeadas
  (Cap.~12), árvores e sistemas grandes usam exatamente esse padrão.
\end{boaprática}
\end{respostabox}

\newpage

% ============================================================
\section{Exercício 10.4 -- Identificar e Corrigir Erros}
% ============================================================

\subsection*{Item (a) -- Precedência de \texttt{*} e \texttt{->}}
\begin{respostabox}
\textbf{Código com erro:}
\begin{lstlisting}
printf( "%s\n", *cPtr->face );
\end{lstlisting}

\textbf{Erro:} O operador \texttt{->} tem precedência \textit{maior} que \texttt{*}.
A expressão é avaliada como \texttt{*(cPtr->face)}, o que tenta desreferenciar o
\texttt{char *} (retornando um \texttt{char}), incompatível com \texttt{\%s} que
espera \texttt{char *}.

\textbf{Correção:} adicionar parênteses para forçar a ordem desejada:
\begin{lstlisting}
printf( "%s\n", ( *cPtr ).face );
\end{lstlisting}
Mais idiomático seria simplesmente:
\begin{lstlisting}
printf( "%s\n", cPtr->face );  /* equivalente */
\end{lstlisting}
\end{respostabox}

\subsection*{Item (b) -- Subscrito de array faltando}
\begin{respostabox}
\textbf{Código com erro:}
\begin{lstlisting}
printf( "%s\n", copas.face );  /* tentava acessar elemento 10 */
\end{lstlisting}

\textbf{Erro:} \texttt{copas} é um \textit{array} (\texttt{hearts[13]}), não uma
estrutura individual. Para acessar o décimo elemento (índice 10), use subscrito:
\begin{lstlisting}
printf( "%s\n", copas[ 10 ].face );
\end{lstlisting}
(Notar que ``elemento 10'' tradicionalmente significa o de índice 9 ou 10 a depender
de convenção; o livro usa 10.)
\end{respostabox}

\subsection*{Item (c) -- Inicialização de \texttt{union}}
\begin{respostabox}
\textbf{Código com erro:}
\begin{lstlisting}
union valores {
    char w;
    float x;
    double y;
};
union valores v = { 1.27 };
\end{lstlisting}

\textbf{Erro:} Em uniões inicializadas com a sintaxe de chaves, a inicialização
aplica-se \textit{apenas ao primeiro membro}. Como o primeiro membro é \texttt{char w},
não dá para inicializá-lo com \texttt{1.27} (um \texttt{double}).

\textbf{Correção:} reordenar a união ou usar inicializador designado (C99):
\begin{lstlisting}
/* Opcao 1: reordenar */
union valores {
    double y;
    char w;
    float x;
};
union valores v = { 1.27 };   /* OK: y = 1.27 */

/* Opcao 2: inicializador designado (C99) */
union valores v = { .y = 1.27 };
\end{lstlisting}
\end{respostabox}

\subsection*{Item (d) -- Falta \texttt{;}}
\begin{respostabox}
\textbf{Código com erro:}
\begin{lstlisting}
struct pessoa {
    char nome[ 15 ];
    char sobrenome[ 15 ];
    int idade;
}
\end{lstlisting}

\textbf{Erro:} Falta o \texttt{;} \textbf{depois} do \texttt{\}}. Diferente do que
acontece com funções (que não têm \texttt{;} após \texttt{\}}), declarações de
\texttt{struct} (e \texttt{union}, \texttt{enum}) precisam de \texttt{;} no final.

\textbf{Correção:}
\begin{lstlisting}
struct pessoa {
    char nome[ 15 ];
    char sobrenome[ 15 ];
    int idade;
};   /* <-- nao esqueca este ; */
\end{lstlisting}

\begin{erroco}
  Esquecer o \texttt{;} após \texttt{\}} de \texttt{struct}/\texttt{union}/\texttt{enum}
  é um erro extremamente comum, e o compilador pode emitir mensagens confusas
  apontando para a próxima linha.
\end{erroco}
\end{respostabox}

\subsection*{Item (e) -- Falta \texttt{struct}}
\begin{respostabox}
\textbf{Código com erro:}
\begin{lstlisting}
pessoa d;   /* assumindo definicao do item (d) ja corrigida */
\end{lstlisting}

\textbf{Erro:} Sem \texttt{typedef}, é necessário usar o nome completo
\texttt{struct pessoa} para declarar variáveis.

\textbf{Correção:}
\begin{lstlisting}
struct pessoa d;
\end{lstlisting}

\textbf{Alternativa:} usar \texttt{typedef}:
\begin{lstlisting}
typedef struct pessoa Pessoa;
Pessoa d;   /* agora funciona */
\end{lstlisting}
\end{respostabox}

\subsection*{Item (f) -- Atribuição entre tipos diferentes}
\begin{respostabox}
\textbf{Código com erro:}
\begin{lstlisting}
struct pessoa p;
struct card   c;
p = c;
\end{lstlisting}

\textbf{Erro:} Em C, atribuição entre \texttt{struct} só funciona entre o
\textbf{mesmo tipo}. \texttt{struct pessoa} e \texttt{struct card} são tipos
distintos, mesmo que tivessem layout idêntico.

\textbf{Não há correção direta:} a operação não faz sentido semanticamente.
Se realmente quiséssemos copiar campos compatíveis, faríamos manualmente:
\begin{lstlisting}
strcpy( p.nome, c.algumCampo );
/* ... outros campos ... */
\end{lstlisting}

\begin{boaprática}
  C é \textbf{tipado nominalmente}: dois tipos \texttt{struct} são diferentes
  mesmo se tiverem o mesmo layout. Isso é uma forma de proteção contra erros
  -- impede misturar conceitos por acidente.
\end{boaprática}
\end{respostabox}

\newpage

% ============================================================
\section*{Gabarito Rápido -- Capítulo 10: Autorrevisão}

\begin{tcolorbox}[colback=deitellightblue,colframe=deitelblue,title={\bfseries\large Respostas Resumidas}]

\textbf{10.1:}
\begin{itemize}[noitemsep]
  \item[a)] estrutura \quad b) união \quad c) AND \texttt{\&}
  \item[d)] membros \quad e) OR inclusivo \texttt{|} \quad f) \texttt{struct}
  \item[g)] \texttt{typedef} \quad h) XOR \texttt{\^{}} \quad i) mascarar
  \item[j)] \texttt{union} \quad k) nome de tag \quad l) \texttt{.} ou \texttt{->}
  \item[m)] \texttt{<<} e \texttt{>>} \quad n) enumeração
\end{itemize}

\medskip\textbf{10.2 -- Verdadeiro/Falso:}
\begin{itemize}[noitemsep]
  \item[a)] F (\texttt{struct} aceita múltiplos tipos)
  \item[b)] F (uniões têm bytes indefinidos)
  \item[c)] V (nome de tag é opcional)
  \item[d)] F (apenas dentro da mesma \texttt{struct})
  \item[e)] F (\texttt{typedef} cria sinônimos, não tipos novos)
  \item[f)] F (\texttt{struct} é passada por valor)
  \item[g)] V (não comparáveis com \texttt{==})
\end{itemize}

\medskip\textbf{10.3:} \texttt{struct peca \{...\}}; \texttt{typedef ... Peca};
\texttt{Peca a, b[10], *ptr;}; \texttt{scanf} com \texttt{\&a.num} e \texttt{a.nome};
\texttt{b[3] = a}; \texttt{ptr = b}; \texttt{(ptr+3)->num}.

\medskip\textbf{10.4 -- Erros corrigidos:}
\begin{itemize}[noitemsep]
  \item[a)] precedência \texttt{*} vs \texttt{->}: use \texttt{(*cPtr).face} ou \texttt{cPtr->face}
  \item[b)] falta subscrito de array
  \item[c)] tipo do primeiro membro da união
  \item[d)] falta \texttt{;} após \texttt{\}} da \texttt{struct}
  \item[e)] falta palavra \texttt{struct} na declaração
  \item[f)] não se pode atribuir entre tipos \texttt{struct} diferentes
\end{itemize}

\end{tcolorbox}

\end{document}
